\chapter{फलन}\label{ch:g10:functions}

\omterm{def:g10:functions:function}{फलन} एक मशीन है जिसमें एक संख्या जाती है और एक संख्या बाहर आती है, और एक ही
निवेश पर हमेशा वही निर्गम मिलता है। आने वाले वर्षों में (\cref{ch:g11:quad} से आगे) आप जो
विश्लेषण खड़ा करेंगे, उसकी केंद्रीय वस्तुएँ \omterm{def:g10:functions:function}{फलन} ही हैं; यह अध्याय शब्दावली
--- \omterm{def:g10:functions:function}{प्रांत}, \omterm{def:g10:functions:function}{प्रतिबिंब}, \omterm{def:g10:functions:graph}{आलेख}, परिवर्तन --- तैयार करता है और \omterm{def:g10:functions:graph}{आलेख} से सूचना
पढ़ने का वह कौशल सधाता है जो आगे अनिवार्य है।

\section{शब्दावली: प्रतिबिंब और पूर्वप्रतिबिंब}

\begin{definition}[फलन]\label{def:g10:functions:function}
मान लीजिए $D$ \omterm{def:g10:numbers:sets}{वास्तविक संख्याओं} का एक समुच्चय है। $D$ पर परिभाषित
\emph{फलन}\index{फलन} $f$ हर संख्या $x \in D$ के साथ ठीक एक \omterm{def:g10:numbers:sets}{वास्तविक संख्या}
जोड़ता है, जिसे $f(x)$ लिखते हैं और $x$ का \emph{प्रतिबिंब}\index{प्रतिबिंब}
कहते हैं। समुच्चय $D$ $f$ का \emph{प्रांत}\index{प्रांत} है। हम लिखते हैं
\[
f : x \longmapsto f(x).
\]
यदि $f(x) = y$ हो, तो $x$ $y$ का \emph{पूर्वप्रतिबिंब}\index{पूर्वप्रतिबिंब}
है।
\end{definition}

\begin{remark}
\omterm{def:g10:functions:function}{प्रांत} के हर $x$ का ठीक एक \omterm{def:g10:functions:function}{प्रतिबिंब} होता है, पर किसी संख्या $y$ के कई
\omterm{def:g10:functions:function}{पूर्वप्रतिबिंब} हो सकते हैं, या एक भी नहीं। $f(x) = x^2$ के लिए: $3$ का \omterm{def:g10:functions:function}{प्रतिबिंब}
$9$ है, और $9$ के दो \omterm{def:g10:functions:function}{पूर्वप्रतिबिंब} $3$ और $-3$ हैं; संख्या $-4$ का कोई
\omterm{def:g10:functions:function}{पूर्वप्रतिबिंब} नहीं है।
\end{remark}

\begin{example}[प्रांत ज्ञात करना]\label{ex:g10:functions:domain}
जब \omterm{def:g10:functions:function}{फलन} किसी सूत्र से दिया जाता है, तब उसका \omterm{def:g10:functions:function}{प्रांत} उन $x$ का समुच्चय होता है
जिनके लिए सूत्र का अर्थ बनता है।
\begin{itemize}
  \item $f(x) = 3x^2 - 5x + 1$ का अर्थ हर $x$ के लिए बनता है: \omterm{def:g10:functions:function}{प्रांत} $\R$ है।
  \item $g(x) = \dfrac{1}{x - 2}$ के लिए $x - 2 \neq 0$ आवश्यक है: \omterm{def:g10:functions:function}{प्रांत} $2$ को छोड़कर सभी \omterm{def:g10:numbers:sets}{वास्तविक
        संख्याएँ} हैं।
  \item $h(x) = \sqrt{x - 3}$ के लिए $x - 3 \geq 0$ आवश्यक है: \omterm{def:g10:functions:function}{प्रांत} $\intco{3}{+\infty}$ है।
\end{itemize}
\end{example}

\begin{example}[प्रतिबिंब और पूर्वप्रतिबिंब निकालना]\label{ex:g10:functions:images}
मान लीजिए $\R$ पर परिभाषित $f(x) = x^2 - 4x + 3$ है।

\emph{$5$ का \omterm{def:g10:functions:function}{प्रतिबिंब}:} $x = 5$ प्रतिस्थापित कीजिए:
$f(5) = 25 - 20 + 3 = 8$।

\emph{$3$ के \omterm{def:g10:functions:function}{पूर्वप्रतिबिंब}:} $f(x) = 3$ हल कीजिए:
\[
x^2 - 4x + 3 = 3
\quad\Longleftrightarrow\quad
x^2 - 4x = 0
\quad\Longleftrightarrow\quad
x(x - 4) = 0,
\]
इसलिए $3$ के \omterm{def:g10:functions:function}{पूर्वप्रतिबिंब} $0$ और $4$ हैं।
\end{example}

\section{फलन का आलेख}

\begin{definition}[आलेख]\label{def:g10:functions:graph}
किसी निर्देशांक-पद्धति में $f$ का \emph{आलेख}\index{आलेख} (या \emph{वक्र})
\omterm{def:g10:functions:function}{प्रांत} के सभी $x$ के लिए बिंदुओं $(x, f(x))$ का समुच्चय है। दूसरे शब्दों में,
बिंदु $(x, y)$ आलेख पर ठीक तब होता है जब $y = f(x)$ हो।
\end{definition}

\begin{omfigure}
\begin{tikzpicture}
\begin{axis}[omaxis,
  width=8.6cm, height=6.0cm,
  xmin=-2.6, xmax=4.6, ymin=-3.4, ymax=4.4,
  xtick={-2,-1,1,2,3,4}, ytick={-3,-2,-1,1,2,3,4}]
  \addplot[omDef, thick, domain=-2.2:4.2, samples=120]
    {0.25*x^3 - 0.75*x^2 - x + 2};
  \draw[dashed, omThm] (axis cs:3,0) -- (axis cs:3,1.25)
    -- (axis cs:0,1.25);
  \fill[omThm] (axis cs:3,1.25) circle (1.5pt);
  \node[omThm, font=\small, above right] at (axis cs:3,1.2)
    {$(3,\,f(3))$};
\end{axis}
\end{tikzpicture}

{\small \omterm{def:g10:functions:graph}{आलेख} पर \omterm{def:g10:functions:function}{प्रतिबिंब} पढ़ना: क्षैतिज अक्ष पर $x = 3$ से चलिए, ऊर्ध्वाधर
दिशा में वक्र तक जाइए, फिर क्षैतिज दिशा में ऊर्ध्वाधर अक्ष तक जाकर $f(3) = 1.25$
पढ़िए।}
\end{omfigure}

\begin{method}[आलेख पढ़ना]\label{met:g10:functions:reading}
$f$ के \omterm{def:g10:functions:graph}{आलेख} पर:
\begin{enumerate}
  \item \emph{$a$ का \omterm{def:g10:functions:function}{प्रतिबिंब}}: क्षैतिज अक्ष पर $x = a$ से ऊर्ध्वाधर दिशा
        में वक्र तक जाइए, फिर क्षैतिज दिशा में ऊर्ध्वाधर अक्ष तक; वहाँ पढ़ा
        गया मान $f(a)$ है;
  \item \emph{$b$ के \omterm{def:g10:functions:function}{पूर्वप्रतिबिंब}}: क्षैतिज रेखा $y = b$ खींचिए;
        \omterm{def:g10:functions:function}{पूर्वप्रतिबिंब} वक्र के साथ उसके सभी प्रतिच्छेद बिंदुओं के $x$-निर्देशांक
        हैं;
  \item \emph{$f(x) = k$ के हल}: वही जो $k$ के \omterm{def:g10:functions:function}{पूर्वप्रतिबिंब} हैं;
  \item \emph{$f(x) \leq k$ के हल}: वे $x$ जिनके लिए वक्र रेखा $y = k$ पर या उसके नीचे
        है।
\end{enumerate}
\end{method}

\begin{omfigure}
\begin{tikzpicture}
\begin{axis}[omaxis,
  width=8.6cm, height=6.0cm,
  xmin=-3.4, xmax=3.4, ymin=-2.4, ymax=4.4,
  xtick={-2,2}, xticklabels={$x_1$,$x_2$}, ytick={2}, yticklabels={$k$}]
  \addplot[omDef, thick, domain=-3:3, samples=120] {0.5*x^2};
  \addplot[omThm, thick, domain=-3.2:3.2] {2};
  \fill (axis cs:-2,2) circle (1.5pt);
  \fill (axis cs:2,2) circle (1.5pt);
  \draw[dashed] (axis cs:-2,2) -- (axis cs:-2,0);
  \draw[dashed] (axis cs:2,2) -- (axis cs:2,0);
\end{axis}
\end{tikzpicture}

{\small $f(x) = k$ को \omterm{def:g10:functions:graph}{आलेख} से हल करना: हल $x_1$ और $x_2$ वक्र के क्षैतिज रेखा $y = k$
के साथ प्रतिच्छेद बिंदुओं के भुज हैं। यहाँ $f(x) \leq k$ $\intcc{x_1}{x_2}$ पर सही है, जहाँ वक्र
रेखा के नीचे है।}
\end{omfigure}

\begin{example}\label{ex:g10:functions:membership}
क्या बिंदु $A(2, 5)$ $f(x) = x^2 + 1$ के \omterm{def:g10:functions:graph}{आलेख} पर है? $f(2) = 4 + 1 = 5$ निकालिए: हाँ, क्योंकि $f(2)$ $A$
के $y$-निर्देशांक के बराबर है। बिंदु $B(3, 8)$ \omterm{def:g10:functions:graph}{आलेख} पर नहीं है, क्योंकि $f(3) = 10 \neq 8$।
\end{example}

\section{फलन के परिवर्तन}

\begin{definition}[वर्धमान, ह्रासमान]\label{def:g10:functions:variations}
मान लीजिए $f$ किसी \omterm{def:g10:numbers:interval}{अंतराल} $I$ पर परिभाषित है।
\begin{itemize}
  \item $f$ $I$ पर \emph{वर्धमान}\index{वर्धमान फलन} है जब सभी $u, v \in I$ के
        लिए, $u < v$ होने पर $f(u) \leq f(v)$ हो: निर्गम निवेश के साथ बढ़ते हैं, और \omterm{def:g10:functions:graph}{आलेख}
        बाएँ से दाएँ चढ़ता है।
  \item $f$ $I$ पर \emph{ह्रासमान}\index{ह्रासमान फलन} है जब $u < v$ से
        $f(u) \geq f(v)$ निकलता हो: \omterm{def:g10:functions:graph}{आलेख} बाएँ से दाएँ उतरता है।
\end{itemize}
कड़ी असमिकाओं $f(u) < f(v)$ (क्रमशः $>$) के साथ हम \emph{निरंतर} वर्धमान (क्रमशः
ह्रासमान) कहते हैं।
\end{definition}

\begin{definition}[परिवर्तन-सारणी]\label{def:g10:functions:table}
\emph{परिवर्तन-सारणी}\index{परिवर्तन-सारणी} यह संक्षेप में बताती है कि $f$
किन \omterm{def:g10:numbers:interval}{अंतरालों} पर बढ़ता और किन पर घटता है --- इसके लिए तीर $\nearrow$ और $\searrow$ का
उपयोग होता है --- और मोड़ के बिंदुओं पर $f$ के मान दर्ज करती है।
\end{definition}

\begin{example}\label{ex:g10:functions:table}
नीचे $\intcc{-2}{4}$ पर आलेखित \omterm{def:g10:functions:function}{फलन} पर विचार कीजिए।

\begin{omfigure}
\begin{tikzpicture}
\begin{axis}[omaxis,
  width=8.6cm, height=5.4cm,
  xmin=-2.8, xmax=4.6, ymin=-2.4, ymax=3.8,
  xtick={-2,1,4}, ytick={-1.5,3}, yticklabels={$-1.5$,$3$}]
  \addplot[omDef, thick, domain=-2:4, samples=120] {-0.5*(x-1)^2 + 3};
  \fill (axis cs:-2,-1.5) circle (1.5pt);
  \fill (axis cs:4,-1.5) circle (1.5pt);
  \fill (axis cs:1,3) circle (1.5pt);
\end{axis}
\end{tikzpicture}

{\small $\intcc{-2}{4}$ पर परिभाषित एक \omterm{def:g10:functions:function}{फलन}, जो पहले बढ़ता है और फिर घटता है।}
\end{omfigure}

\omterm{def:g10:functions:graph}{आलेख} पढ़ने पर: $f$ $f(-2) = -1.5$ से बढ़कर अपने \omterm{def:g10:functions:extrema}{अधिकतम} $f(1) = 3$ तक जाता है, फिर घटकर
$f(4) = -1.5$ तक आता है। इसकी \omterm{def:g10:functions:table}{परिवर्तन-सारणी} है
\begin{center}
\renewcommand{\arraystretch}{1.3}
\begin{tabular}{c|ccccc}
$x$ & $-2$ & & $1$ & & $4$ \\
\hline
$f$ & $-1.5$ & $\nearrow$ & $3$ & $\searrow$ & $-1.5$ \\
\end{tabular}
\end{center}
\end{example}

\begin{example}[परिवर्तन सिद्ध करना]\label{ex:g10:functions:proof}
दिखाइए कि $f(x) = 3x + 1$ $\R$ पर निरंतर \omterm{def:g10:functions:variations}{वर्धमान} है। कोई भी $u < v$ लीजिए और
\omterm{def:g10:functions:function}{प्रतिबिंबों} की तुलना कीजिए:
\[
f(v) - f(u) = (3v + 1) - (3u + 1) = 3(v - u) > 0,
\]
क्योंकि $v - u > 0$ है। तो $f(u) < f(v)$: \omterm{def:g10:functions:function}{फलन} निरंतर \omterm{def:g10:functions:variations}{वर्धमान} है। ऋणात्मक प्रवणता के साथ
वही गणना, जैसे $g(x) = -2x + 5$, $g(v) - g(u) = -2(v-u) < 0$ देती है: निरंतर \omterm{def:g10:functions:variations}{ह्रासमान}।
\end{example}

\section{चरम मान}

\begin{definition}[अधिकतम, न्यूनतम]\label{def:g10:functions:extrema}
मान लीजिए $f$ $D$ पर परिभाषित है और $a \in D$ है। मान $f(a)$ $D$ पर $f$
का \emph{अधिकतम}\index{अधिकतम} है जब सभी $x \in D$ के लिए $f(x) \leq f(a)$ हो; और वह
\emph{न्यूनतम}\index{न्यूनतम} है जब सभी $x \in D$ के लिए $f(x) \geq f(a)$ हो। \omterm{def:g10:functions:graph}{आलेख} पर ये
वक्र के सबसे ऊँचे और सबसे नीचे बिंदु हैं।
\end{definition}

\begin{example}\label{ex:g10:functions:extremum}
दिखाइए कि $f(x) = (x - 1)^2 + 2$ का $\R$ पर \omterm{def:g10:functions:extrema}{न्यूनतम} $2$ है, जो $x = 1$ पर प्राप्त होता है।
हर $x$ के लिए वर्ग $(x-1)^2$ $\geq 0$ है, इसलिए $f(x) = (x-1)^2 + 2 \geq 2$; और $f(1) = 0 + 2 = 2$, इसलिए मान $2$
सचमुच प्राप्त होता है। ये दोनों तथ्य मिलकर सिद्ध करते हैं कि \omterm{def:g10:functions:extrema}{न्यूनतम} $2$
है।
\end{example}

\begin{omfigure}
\begin{tikzpicture}
\begin{axis}[omaxis,
  width=7.6cm, height=5.2cm,
  xmin=-1.8, xmax=3.8, ymin=-0.9, ymax=6.4,
  xtick={1}, ytick={2}]
  \addplot[omDef, thick, domain=-1.0:3.0, samples=100] {(x-1)^2 + 2};
  \addplot[omThm, dashed, domain=-1.6:3.6] {2};
  \fill (axis cs:1,2) circle (1.5pt);
  \node[below right, font=\small] at (axis cs:1,2) {$(1,2)$};
\end{axis}
\end{tikzpicture}

{\small $f(x) = (x-1)^2 + 2$ का वक्र बिंदुदार रेखा $y = 2$ से कभी नीचे नहीं जाता, और $x = 1$
पर उसे छूता है: $f$ का \omterm{def:g10:functions:extrema}{न्यूनतम} $2$ है।}
\end{omfigure}

\section{अभ्यास}

\begin{exercise}[$\star$]\label{exo:g10:functions:1}
मान लीजिए $f(x) = 2x^2 - 3x + 1$ है। $0$, $2$, $-1$ और $\frac12$ के \omterm{def:g10:functions:function}{प्रतिबिंब} निकालिए।
\end{exercise}

\begin{exercise}[$\star$]\label{exo:g10:functions:2}
हर \omterm{def:g10:functions:function}{फलन} का \omterm{def:g10:functions:function}{प्रांत} बताइए:
\[
f(x) = 5x - 2, \qquad
g(x) = \frac{3}{x + 4}, \qquad
h(x) = \sqrt{2x - 6}, \qquad
k(x) = \frac{1}{x^2 + 1}.
\]
\end{exercise}

\begin{exercise}[$\star$]\label{exo:g10:functions:3}
मान लीजिए $f(x) = x^2 - 2x$ है। $0$, $3$ और $-1$ के सभी \omterm{def:g10:functions:function}{पूर्वप्रतिबिंब} ज्ञात
कीजिए।
\end{exercise}

\begin{exercise}[$\star$]\label{exo:g10:functions:4}
क्या बिंदु $A(-1, 4)$ $f(x) = x^2 - 2x + 1$ के \omterm{def:g10:functions:graph}{आलेख} का है? और बिंदु $B(2, 1)$? गणना से पुष्ट कीजिए।
\end{exercise}

\begin{exercise}[$\star$]\label{exo:g10:functions:5}
$\intcc{-3}{5}$ पर परिभाषित \omterm{def:g10:functions:function}{फलन} $g$ की \omterm{def:g10:functions:table}{परिवर्तन-सारणी} यह है
\begin{center}
\renewcommand{\arraystretch}{1.3}
\begin{tabular}{c|ccccccc}
$x$ & $-3$ & & $0$ & & $3$ & & $5$ \\
\hline
$g$ & $1$ & $\searrow$ & $-2$ & $\nearrow$ & $4$ & $\searrow$ & $0$ \\
\end{tabular}
\end{center}
\begin{enumerate}
  \item $\intcc{-3}{5}$ पर $g$ के \omterm{def:g10:functions:extrema}{अधिकतम} और \omterm{def:g10:functions:extrema}{न्यूनतम} क्या हैं?
  \item $g(-1)$ और $g(-0.5)$ के मान जाने बिना उनकी तुलना कीजिए।
  \item सारणी के हर \omterm{def:g10:numbers:interval}{अंतराल} पर \omterm{def:g10:algebra:equation}{समीकरण} $g(x) = 0$ के अधिक से अधिक कितने हल हो
        सकते हैं?
\end{enumerate}
\end{exercise}

\begin{exercise}[$\star\star$]\label{exo:g10:functions:6}
मान लीजिए $f(x) = -2x + 7$ है। $u < v$ के लिए $f(u)$ और $f(v)$ की तुलना करके दिखाइए कि
$f$ $\R$ पर निरंतर \omterm{def:g10:functions:variations}{ह्रासमान} है।
\end{exercise}

\begin{exercise}[$\star\star$]\label{exo:g10:functions:7}
दिखाइए कि \omterm{def:g10:functions:function}{फलन} $f(x) = x^2 + 6x + 11$ का $\R$ पर \omterm{def:g10:functions:extrema}{न्यूनतम} है, और उसका मान तथा वह स्थान
बताइए जहाँ वह प्राप्त होता है। (संकेत: उपयुक्त अचर $c$ लेकर $f(x) = (x + 3)^2 + c$
लिखिए।)
\end{exercise}

\begin{exercise}[$\star\star$]\label{exo:g10:functions:8}
एक आयताकार बग़ीचे का परिमाप $40$ m है। मान लीजिए $x$ मीटर में उसकी
चौड़ाई है।
\begin{enumerate}
  \item लंबाई को, और फिर क्षेत्रफल $A(x)$ को, $x$ के \omterm{def:g10:functions:function}{फलन} के रूप में
        व्यक्त कीजिए। किन $x$ के लिए इसका अर्थ बनता है?
  \item $A(4)$, $A(8)$, $A(10)$, $A(12)$ और $A(16)$ निकालिए। सबसे बड़े क्षेत्रफल
        वाली आकृति के बारे में आपका अनुमान क्या है?
\end{enumerate}
\end{exercise}

\begin{exercise}[$\star\star$]\label{exo:g10:functions:9}
मान लीजिए $\R$ पर परिभाषित $f(x) = \dfrac{1}{x^2 + 1}$ है।
\begin{enumerate}
  \item दिखाइए कि हर वास्तविक $x$ के लिए $0 < f(x) \leq 1$ है।
  \item क्या $\R$ पर $f$ का कोई \omterm{def:g10:functions:extrema}{अधिकतम} है? कोई \omterm{def:g10:functions:extrema}{न्यूनतम}? पुष्ट कीजिए।
\end{enumerate}
\end{exercise}

\begin{exercise}[$\star\star\star$]\label{exo:g10:functions:10}
मान लीजिए $\R$ पर $f(x) = x^2$ है।
\begin{enumerate}
  \item मान लीजिए $0 \leq u < v$ है। $f(v) - f(u)$ का \omterm{def:g10:algebra:expand}{गुणनखंडन} कीजिए और निष्कर्ष निकालिए
        कि $f$ $\intco{0}{+\infty}$ पर निरंतर \omterm{def:g10:functions:variations}{वर्धमान} है।
  \item यही तर्क ढालकर दिखाइए कि $f$ $\intoc{-\infty}{0}$ पर निरंतर \omterm{def:g10:functions:variations}{ह्रासमान} है।
\end{enumerate}
\end{exercise}

\section{समस्या: अपना ही निर्गम खाने वाली मशीनें}

\begin{problem}[{सप्ताहांत समस्या --- स्थिर बिंदु, पुनरावृत्ति, और तीन
प्रसिद्ध मशीनें: हीरोन की सुधारक, वर्गकारी, और अनसुलझी ओला-मशीन}]
\label{pb:g10:functions:1}
\omterm{def:g10:functions:function}{फलन} एक मशीन है: एक संख्या भीतर जाती है, एक संख्या बाहर आती है। सबसे रोचक
प्रयोग वे हैं जिनमें मशीन को \emph{उसी का निर्गम} बार-बार खिलाया जाता है
--- और तब दो प्रश्न सब पर हावी हो जाते हैं: क्या प्रक्रिया कहीं जाकर ठहरती
है (कोई \emph{\omterm{pb:g10:functions:1}{स्थिर बिंदु}}), और मशीन का परिवर्तन उसे वहाँ तक कैसे ले जाता
है? इस समस्या की एक मशीन दो हज़ार साल से वर्गमूलों को माँज रही है; दूसरी
प्रारंभिक गणित की सबसे प्रसिद्ध अनसुलझी समस्या की पहरेदार है।

\medskip
\noindent\textbf{भाग I --- हीरोन की मशीन।}
इस \omterm{def:g10:functions:function}{फलन} पर विचार कीजिए:
\[
f(x) = \frac12\left(x + \frac2x\right).
\]
\begin{enumerate}
  \item $f$ का \omterm{def:g10:functions:function}{प्रांत} बताइए, और $f(1)$, $f(2)$ तथा $f\!\left(\frac32\right)$ को यथार्थ भिन्नों
        के रूप में निकालिए।
  \item $\frac32$ के सभी \omterm{def:g10:functions:function}{पूर्वप्रतिबिंब} ज्ञात कीजिए: $f(x) = \frac32$ हल कीजिए (हर हटाइए
        और \cref{pb:g10:algebra:1} की विधियाँ काम में लाइए)।
  \item $f$ का \emph{स्थिर बिंदु}\index{स्थिर बिंदु} वह संख्या है जिसे
        मशीन बदलती नहीं: $f(x) = x$। $f$ के दोनों \omterm{pb:g10:functions:1}{स्थिर बिंदु} ज्ञात कीजिए।
        धनात्मक वाले की पहरेदारी कौन-सा पुराना परिचित करता है?
  \item $x = 1$ से शुरू करके मशीन को उसका अपना निर्गम तीन बार खिलाइए और
        यथार्थ परिणाम दर्ज कीजिए। माध्यमिक विद्यालय खंड की अपरिमेयता वाली
        सप्ताहांत समस्या का कौन-सा अनुक्रम आपने फिर से खड़ा कर दिया है ---
        और \omterm{def:g10:functions:function}{फलनों} की भाषा में ``हीरोन का नुस्ख़ा'' क्या था?
  \item एक ही चित्र में $f$ का \omterm{def:g10:functions:graph}{आलेख} ($x > 0$ के लिए) और विकर्ण रेखा $y = x$
        खींचिए। इस चित्र में \omterm{pb:g10:functions:1}{स्थिर बिंदु} कहाँ हैं? (माध्यमिक विद्यालय खंड
        की दो-थर्मामीटर वाली सप्ताहांत समस्या से तुलना कीजिए: दोनों
        मापक्रमों पर एक ही ताप पढ़ना यही विचार था।)
\end{enumerate}

\medskip
\noindent\textbf{भाग II --- मशीन चूक क्यों नहीं सकती।}
\begin{enumerate}[resume]
  \item $0 < u < v$ के लिए उभयनिष्ठ हर पर लाकर दिखाइए कि
        \[
        f(v) - f(u) = \frac{(v - u)(uv - 2)}{2uv} ,
        \]
        और $\intoo{0}{+\infty}$ पर $f$ के परिवर्तन निकालिए: एक बिंदु तक घटता, फिर बढ़ता
        --- कौन-सा बिंदु? (\cref{exo:g10:functions:10} वाली वही तकनीक।)
  \item निष्कर्ष निकालिए कि $\intoo{0}{+\infty}$ पर \omterm{def:g10:functions:function}{फलन} $f$ का एक \omterm{def:g10:functions:extrema}{न्यूनतम} है, जिसका मान
        ठीक $\sqrt2$ है और जो $\sqrt2$ पर प्राप्त होता है। मशीन के \emph{हर}
        निर्गम (धनात्मक निवेशों के लिए) के बारे में यह क्या कहता है?
  \item दिखाइए कि यदि $x > \sqrt2$ हो तो $\sqrt2 < f(x) < x$ होता है: बहुत बड़ा अनुमान हमेशा
        सुधरता है, कभी उससे नीचे नहीं गिरता। ($f(x) < x$ के लिए $\frac2x$ की $x$
        से तुलना कीजिए।)
  \item प्रश्न 6--8 को जोड़कर $\intoo{0}{+\infty}$ पर $f$ की \omterm{def:g10:functions:table}{परिवर्तन-सारणी} बनाइए
        (\cref{def:g10:functions:table}), \omterm{def:g10:functions:extrema}{न्यूनतम} सहित।
  \item दूसरी मशीन: $g(x) = x^2$। इसके \omterm{pb:g10:functions:1}{स्थिर बिंदु} $0$ और $1$ हैं (जाँचिए)।
        $g$ को $x = 0.9$ से चार बार दोहराइए, फिर $x = 1.1$ से चार बार
        (कैलकुलेटर, तीन दशमलव स्थान)। एक \omterm{pb:g10:functions:1}{स्थिर बिंदु} खींचता है, दूसरा
        धकेलता है: कौन-सा कौन है?
\end{enumerate}

\medskip
\noindent\textbf{भाग III --- ओला-मशीन।}
पूर्ण संख्याओं पर परिभाषित कीजिए: यदि $n$ सम हो तो $h(n) = \frac n2$, और यदि $n$
विषम हो तो $h(n) = 3n + 1$। $h$ के अधीन संख्याएँ तूफ़ानी बादल में ओलों की तरह उछलती
हैं।
\begin{enumerate}[resume]
  \item $h$ के अधीन $7$ की पूरी उड़ान $1$ तक निकालिए। इसमें कितने
        चरण लगते हैं, और वह कितनी ऊँची उड़ती है?
  \item $27$ की उड़ान शुरू कीजिए और दस चरण निकालिए। (इसकी पूरी उड़ान $111$
        चरणों तक चलती है और $9\,232$ पर चोटी छूती है --- और शुरुआत $27$ से
        होती है।) सरल नियमों के बारे में $7$ और $27$ क्या सिखाते हैं?
  \item समझाइए कि $2$ की किसी घात तक पहुँचने वाली कोई भी उड़ान सीधे $1$
        तक क्यों गिर जाती है, और $1$ पर क्या होता है ($h(1)$, $h(4)$, $h(2)$
        निकालिए)। क्या धनात्मक पूर्ण संख्याओं में $h$ का कोई \omterm{pb:g10:functions:1}{स्थिर बिंदु}
        है? इसके बदले हर देखी गई उड़ान का अंत क्या बनाता है?
  \item \emph{कोलात्ज़ अनुमान} कहता है कि हर आरंभिक संख्या अंततः $1$ तक
        पहुँच जाती है। संगणकों ने इसे $10^{20}$ से बहुत आगे तक जाँच लिया है;
        किसी मनुष्य ने इसे सिद्ध नहीं किया। किन दो प्रकार की खोजों से यह
        \emph{ग़लत} सिद्ध होगा? और भारी-भरकम जाँच के बाद भी मामला बंद क्यों
        नहीं होता (माध्यमिक विद्यालय खंड की माचिस-भविष्यवक्ता वाली
        सप्ताहांत समस्या ने यही घंटी बजाई थी)?
  \item ``गणित अभी ऐसे प्रश्नों के लिए पका नहीं है,'' पॉल एर्दॉश ने इस
        अनुमान के बारे में कहा था। एक वाक्य में: ओला-मशीन को हीरोन की मशीन
        से क्या अलग करता है, जहाँ प्रश्न 6--8 ने सब कुछ तय कर दिया था?
\end{enumerate}

\medskip
\noindent\textbf{भाग IV --- मशीन चलाने की निर्देश-पुस्तिका।}
\begin{enumerate}[resume]
  \item \omterm{def:g10:numbers:interval}{अंतरालों} के रूप में \omterm{def:g10:functions:function}{प्रांत} (\cref{pb:g10:numbers:1} की भाषा): $x \mapsto \sqrt{x - 3}$, \ $x \mapsto \frac{1}{x^2 - 9}$, \
        $x \mapsto \sqrt{9 - x^2}$ के \omterm{def:g10:functions:function}{प्रांत} बताइए।
  \item हीरोन के $f$ पर लौटिए: $3$ के यथार्थ \omterm{def:g10:functions:function}{पूर्वप्रतिबिंब} ज्ञात
        कीजिए ($f(x) = 3$ को वर्ग पूरा करके हल कीजिए)।
  \item एक पत्थर ऊपर की ओर फेंका जाता है: $t$ सेकंड बाद उसकी ऊँचाई $H(t) = 20t - 5t^2$
        मीटर होती है। किस समय-अंतराल पर इस सूत्र का भौतिक अर्थ बनता है?
        वर्ग पूरा करके \omterm{def:g10:functions:extrema}{अधिकतम} ऊँचाई और उसका समय ज्ञात कीजिए (\cref{def:g10:functions:extrema}), और
        \omterm{def:g10:functions:table}{परिवर्तन-सारणी} दीजिए।
  \item पत्थर कितनी देर तक कम से कम $15$ m ऊँचाई पर रहता है? ($H(t) \geq 15$ को
        \omterm{def:g10:algebra:signtable}{चिह्न-सारणी} से हल कीजिए, \cref{met:g10:algebra:signtable}।)
  \item समापन: इस समस्या ने चार मशीनें चलाईं --- हीरोन की सुधारक,
        वर्गकारी, ओला-मशीन, और पत्थर की ऊँचाई। दो-तीन वाक्यों में बताइए कि
        हर \omterm{def:g10:functions:function}{फलन} के पास क्या-क्या होता है (\omterm{def:g10:functions:function}{प्रांत}, नियम, \omterm{def:g10:functions:graph}{आलेख}) और इस समस्या
        ने हर मशीन से कौन-से दो बड़े प्रश्न पूछे (पुनरावृत्ति कहाँ जाकर
        ठहरती है? \omterm{def:g10:functions:function}{फलन} कैसे बदलता है?) --- वही प्रश्न आगे के अध्याय एक-एक
        करके उठाते हैं।
\end{enumerate}
\end{problem}
